\SourceLeaf{253}
\par\bigskip\Needspace{9\baselineskip}
\pdfbookmark[0]{Disquisitionum circa aequationes puras ulterior evolutio}{aequationes-purae}
\begin{center}\large\bfseries DISQUISITIONUM CIRCA AEQUATIONES PURAS\par\medskip ULTERIOR EVOLUTIO.\end{center}
\BookRule{.12\linewidth}
\renewcommand{\BookRunningTitle}{Nachlass. Disquisitionum circa aequationes puras}
\renewcommand{\Article}[1]{\par\medskip\Needspace{4\baselineskip}\pdfbookmark[1]{Art. #1}{purae-art-#1}\centerline{\textbf{#1.}}\nobreak\smallskip}

\Article{1}
Quum methodus ea, per quam in \emph{Disquiss.\ Arithm.}\ art.\ 360 aequationem $x^n-1=0$ solvere docuimus, theoriam foecundissimam et gravissimam constituat, cuius prima tantum momenta in opere illo attingere licuit, gratum geometris fore speramus, si hoc argumentum denuo hic resumimus, quae breviter tantum partimque demonstrationibus suppressis adumbrata fuerant, uberius tractamus, et quae ex illo tempore accesserunt incrementa profundius persequimur.

Exponens $n$ supponitur esse numerus primus, numerusque $n-1$ in factores $\alpha\times\beta\times\gamma$ resolutus; porro designamus per $g$ aliquam radicem primitivam pro modulo $n$. Exhibeat $r$ indefinite radicem aequationis $x^n-1=0$, atque $R$ indefinite radicem aequationis $x^\beta-1=0$. Designando itaque peripheriam circuli, cuius radius $=1$, per $P$, quantitatemque imaginariam $\mathord{\surd}{-1}$ per $i$, omnes radices aequationis $x^\beta-1=0$, sive omnes valores ipsius $R$ exhibebuntur per formulam
\[
\cos\frac{kP}{\theta}+i\sin\frac{kP}{\theta}
\]
exprimente $k$ indefinite numeros integros $0,1,2,3\ldots\beta-1$. Porro patet, omnes potestates cuiusvis radicis $R$ ipsas quoque esse radices, nec non, si $R$ fuerit radix valori ipsius $k$ ad $\beta$ primo respondens, omnes potestates $R^0,R,R^2,R^3\ldots R^{\beta-1}$ inter se diversas esse, adeoque totum radicum complexum exhaurire; in hoc casu ipsam $R$ radicem \emph{propriam} aequationis $x^\beta-1=0$ dicemus; contra radix $R$ 
\SourceLeaf{254}
valori ipsius $k$ ad $\beta$ non primo respondens \emph{impropria} vocabitur, nulloque negotio perspicitur, si $\delta$ fuerit divisor communis maximus numerorum $k$ et $\beta$, fore $R^{\frac\beta\delta .}=1$, omnes vero potestates $R^0,R,R^2,R^3\ldots R^{\frac\beta\delta-1}$ inter se diversas, adeoque $R$ radicem propriam aequationis $x^{\frac\beta\delta}-1=0$. Eadem de aequatione $x^n-1=0$ valebunt, sed huius radices omnes necessario sunt propriae radice $1$ excepta.

\Article{2}
His praemissis disquisitio nostra imprimis versabitur circa functiones huius formae, e $\beta\gamma$ terminis conflatas
\[
r+Rr^{g^\alpha}+R^2r^{g^{2\alpha}}+R^3r^{g^{3\alpha}}\ldots+R^{\beta\gamma-1}r^{g^{\alpha\beta\gamma-\alpha}}
\]
quas compendii caussa per hunc characterem $[r,R]$ designabimus. Singuli termini talis expressionis sunt producta e potestatibus ipsius $r$ in potestates ipsius $R$; illarum exponentes progressionem geometricam constituunt, exponentes harum arithmeticam. Exponentes
\[
1,\quad g^\alpha,\quad g^{2\alpha},\quad g^{3\alpha}\ldots\ldots g^{\alpha\beta\gamma-\alpha}
\]
omnes inter se incongrui sunt secundum modulum $n$, adeoque illae potestates ipsius $r$ inter se diversae; ulterius vero continuatae eandem seriem denuo inciperent, quum sit $g^{\alpha\beta\gamma}\equiv1(\mathrm{mod.}\,n)$, adeoque $r^{g^{\alpha\beta\gamma}}=r$. Factores alteri autem
\[
1,\quad R,\quad R^2,\quad R^3\dotsfour R^{\beta\gamma-1}
\]
constituunt $\gamma$ periodos aequales, quum sit $R^\beta=1$, $R^{\beta-1}=R$ etc. Hinc patet, functionem $[r,R]$ ita quoque exhiberi posse
\[
\begin{aligned}
&r+r^{g^{\alpha\beta}}+r^{g^{2\alpha\beta}}\ldots\ldots+r^{g^{\alpha\beta\gamma-\alpha\beta}}\\
&{}+R\bigl(r^{g^\alpha}+r^{g^{\alpha\beta+\alpha}}+r^{g^{2\alpha\beta+\alpha}}\ldots\ldots+r^{g^{\alpha\beta\gamma-\alpha\beta+\alpha}}\bigr)\\
&{}+R^2\bigl(r^{g^{2\alpha}}+r^{g^{\alpha\beta+2\alpha}}+r^{g^{2\alpha\beta+2\alpha}}\ldots\ldots+r^{g^{\alpha\beta\gamma-\alpha\beta+2\alpha}}\bigr)\\
&{}+\text{etc.}\\
&{}+R^{\beta-1}\bigl(r^{g^{\alpha\beta-\alpha}}+r^{g^{2\alpha\beta-\alpha}}+r^{g^{3\alpha\beta-\alpha}}\ldots\ldots+r^{g^{\alpha\beta\gamma-\alpha}}\bigr).
\end{aligned}
\]
sive introducendo signum art.\ 343 Disq.\ Ar.
\[
[r,R]=(\gamma,1)+R(\gamma,g^\alpha)+R^2(\gamma,g^{2\alpha})\ldots+R^{\beta-1}(\gamma,g^{\alpha\beta-\alpha})
\]

\SourceLeaf{255}
\Article{3}
Si pro radice $r$ unitatem accipimus, habemus
\[
[1,R]=1+R+R^2+R^3\ldots+R^{\beta\gamma-1}=\gamma(1+R+R^2+R^3\ldots+R^{\beta-1})
\]
huius valor erit $=\beta\gamma$, si etiam pro $R$ accipitur radix $1$, sed $=0$ pro quovis alio valore ipsius $R$. Contra manente $r$ indeterminata, positaque $R=1$, erit $[r,1]=r+r^{g^\alpha}+r^{g^{2\alpha}}+r^{g^{3\alpha}}\ldots+r^{g^{\alpha\beta\gamma-\alpha}}$, sive adhibito signo in Disq.\ Ar.\ introducto, $[r,1]=(\beta\gamma,1)$, i.\ e.\ constabit e periodo $\beta\gamma$ radicum, e quibus una est ipsa $r$. Quoties est $\alpha=1$, haec periodus omnes radices $r,r^2,r^3\ldots r^{n-1}$ complectetur ordine tantum mutato.

Notentur adhuc relationes sequentes, quarum ratio sponte elucet:
\[
[r,R]=R[r^{g^\alpha},R]=R^2[r^{g^{2\alpha}},R]\quad\text{sive generaliter}\quad=R^k[r^{g^{\alpha k}},R]
\]
denotante $k$ integrum positivum quemcunque. Hinc patet, functionem $[r^m,R]$ vel esse $=[1,R]$, scilicet si fuerit $m$ divisibilis per $n$, vel reduci posse ad formam $R^\mu[r^{g^\nu},R]$ in casibus reliquis et quidem ita, ut sit $\nu<\alpha$. Si enim $m$ non est divisibilis per $n$, congruus erit secundum modulum $n$ alicui potestati ipsius $g$, cuius exponens ad instar Disq.\ Ar.\ per $\mathrm{ind.}\,m$ commode exprimitur; statuendo itaque $\mathrm{ind.}\,m=\lambda\alpha+\nu$, quod manifesto fieri potest, ita ut sit $\nu<\alpha$, erit $[r^m,R]=[r^{g^{\lambda\alpha+\nu}},R]=R^{-\lambda}[r^{g^\nu},R]$: faciendus est itaque $\mu=-\lambda$ aut si exponentem positivum desideras, $\mu\equiv-\lambda(\mathrm{mod.}\,\beta)$.

\Article{4}
\textsc{Theorema.} \emph{Designante $r'$ perinde ut $r$ indefinite radicem aequationis $x^n-1=0$, nec non $R'$ perinde ut $R$ indefinite radicem aequationis $x^\beta-1=0$, erit productum}
\[
\begin{aligned}
&[r,R]\times[r',R']=\\
&\quad[rr',RR']+R[r^{g^\alpha}r',RR']+R^2[r^{g^{2\alpha}}r',RR']\\
&\quad{}+R^3[r^{g^{3\alpha}}r',RR']\ldots+R^{\beta\gamma-1}[r^{g^{\alpha\beta\gamma-\alpha}}r',RR']
\end{aligned}
\]

\emph{Demonstr.} Absolvendo multiplicationem ipsius $[r,R]$ per singulas partes ipsius $[r',R']$, productum in hac forma exhiberi potest
\[
\begin{aligned}
&[r,R]r'+RR'[r^{g^\alpha},R]r'^{g^\alpha}+R^2R'^2[r^{g^{2\alpha}},R]r'^{g^{2\alpha}}\\
&{}+R^3R'^3[r^{g^{3\alpha}},R]r'^{g^{3\alpha}}\ldots\\
&{}+R^{\beta\gamma-1}R'^{\beta\gamma-1}[r^{g^{\alpha\beta\gamma-\alpha}},R]r'^{g^{\alpha\beta\gamma-\alpha}}
\end{aligned}
\]

\SourceLeaf{256}
Collectis dein singularum partium rite evolutarum terminis primis, prodit $[rr',RR']$; perinde collectis terminis secundis, emergit $R[r^{g^\alpha}r',RR']$ et sic porro, unde tandem producti forma tradita conflatur. Q.\ E.\ D.

Ceterum per solam permutationem ipsarum $r,R$ cum $r',R'$ patet, idem productum etiam sub hanc formam poni posse:
\[
\begin{aligned}
&[rr',RR']+R'[rr'^{g^\alpha},RR']+R'^2[rr'^{g^{2\alpha}},RR']\\
&{}+R'^3[rr'^{g^{3\alpha}},RR']\dotsfour+R'^{\beta\gamma-1}[rr'^{g^{\alpha\beta\gamma-\alpha}},RR']
\end{aligned}
\]

Hinc porro concluditur, si etiam $r'',r'''$ etc.\ indefinite exprimant radices aequationis $x^n-1=0$, nec non $R'',R'''$ etc.\ indefinite radices aequationis $x^\beta-1=0$, productum e functionibus $[r,R],[r',R'],[r'',R''],[r''',R''']$ etc., quantacunque fuerit ipsarum multitudo, aequale fore aggregato
\[
\begin{gathered}
\Sigma R'^{k'}R''^{k''}R'''^{k'''}\text{ etc.}\\
{}[rr'^{g^{\alpha k'}}r''^{g^{\alpha k''}}r'''^{g^{\alpha k'''}}\text{ etc.},\ RR'R''R'''\text{ etc.}]
\end{gathered}
\]
substitutis pro $k',k'',k'''$ etc.\ omnibus numeris $0,1,2,3\ldots\beta\gamma-1$, omnibus modis diversis possibilibus inter se combinatis, quo pacto omnino $\beta^{\mu-1}\gamma^{\mu-1}$ termini emergent, si per $\mu$ multitudo illarum functionum inter se multiplicatarum denotatur.
